% Revised project proposal (Group 6) -- addresses the feedback of 8 Oct 2026:
% (1) entity = edition, (2) complete DBpedia download with dbp: properties,
% (3) documented Amazon selection with random baseline.
\documentclass[runningheads]{llncs}
\usepackage[T1]{fontenc}
\usepackage{graphicx}
\usepackage{tabularx}
\usepackage{xurl}
\usepackage{booktabs}
\usepackage{array}
\usepackage{threeparttable}
\usepackage{ragged2e}
\setlength{\textfloatsep}{8pt plus 2pt minus 2pt}
\setlength{\intextsep}{8pt plus 2pt minus 2pt}
\setlength{\floatsep}{6pt plus 2pt minus 2pt}
\setlength{\emergencystretch}{2em}

\begin{document}
\title{Web Data Integration Project Proposal (Group 6) -- Revised}
\titlerunning{Book Data Integration}
\authorrunning{Bolduan et al.}
\author{Lucia Bolduan \orcidID{2261625} \and Safa Shaghely \orcidID{2480414} \and Selma Nezihoglu \orcidID{2390917} \and Minji Oh \orcidID{2389125} \and Kejsi Limaci \orcidID{2389812} \and Shiqi Xu \orcidID{1863435}}
\institute{University of Mannheim}
\maketitle

\begin{abstract}
We integrate book data from Goodreads (52,478 records), DBpedia (59,270) and
Amazon (40,000) into a common schema. Following the feedback, the entity is a
book \emph{edition} (e.g.\ the 2018 Putnam hardcover of \emph{Where the Crawdads
Sing}), not a work. Using the ISBN only as an evaluation key, 6,087 editions
occur in at least two datasets. Many works appear in several editions with the
same title and author (6,588 works with 33,540 editions in the Amazon subset),
which makes telling editions apart the central challenge for identity
resolution. The integrated schema has 14 attributes, 12 of which occur in at
least two datasets; \texttt{authors} and \texttt{genres} are list attributes.
\end{abstract}

\section{Use Case and Entity Definition}

The goal is to integrate book information from three independently maintained
sources that describe overlapping books with different attribute names,
formats and levels of detail. One record of the integrated dataset describes
one \textbf{edition}: a specific publication of a work by a publisher, in a
format (hardcover, paperback, \dots), at a date, with its own page count and ISBN.
Goodreads and Amazon list editions directly. A DBpedia resource describes a
work, but its infobox values (ISBN, publisher, release date, pages) are those
of the first edition, so each DBpedia record is treated as the first edition of
its work. Matching therefore has to distinguish, e.g., the hardcover, paperback,
large-print and movie tie-in editions of the same novel, which share title and
author but differ in publisher, date, format and page count.

\section{Datasets}

\textbf{Goodreads.} The ``Best Books Ever'' dataset (scraped from the Goodreads
list of the same name) lists one edition per book. We use the original CSV;
an earlier copy had been re-saved with a spreadsheet program, which had turned
91\% of the ISBNs into scientific notation and altered dates.

\textbf{DBpedia.} All 59,270 resources of type \texttt{dbo:Book} were
downloaded from the public SPARQL endpoint. Because the endpoint silently cuts
off long-running queries, the download (1) lists the book URIs in partitions by
first character, paged with \texttt{LIMIT 10000 OFFSET} (below the 40,000 offset
limit), and (2) fetches the attributes for 150 known URIs per query
(\texttt{VALUES}); a block that returns a partial result (HTTP 206) is split
and fetched again. The final count equals \texttt{COUNT(dbo:Book)}. Year and
language combine \texttt{dbo:releaseDate}/\texttt{dbo:language} with the raw
infobox properties \texttt{dbp:releaseDate}, \texttt{dbp:pubDate} and
\texttt{dbp:language}, which lowers missing values from 88.5\% to 11.4\%
(year) and from 87.6\% to 20.4\% (language). Unlike the first version, books
without an ISBN are no longer excluded, so 39.0\% of the records have no ISBN.

\textbf{Amazon.} The metadata file of Amazon Reviews 2023 (McAuley Lab)
contains 4,448,181 book products. 4,336,308 have a title and an author, and
3,526,297 of these also have an ISBN; this is the population we sample from.
Products without ISBN (mostly Kindle and Audible) are excluded because the
ISBN is our gold-standard key. From this population we take the
\textbf{40,000 editions with the most ratings} (at least 6,506 ratings each).
The selection uses only Amazon's own rating counts, no information from the
other datasets. We chose popularity because Goodreads ``Best Books Ever'' is
itself a popularity-based list: a uniform random sample of 40,000 editions
(fixed seed) shares only 331 editions with Goodreads, compared with 2,169 for
the popularity-based subset (Table~\ref{tab:overlap}). Amazon reports one
rating count for all editions of a work, so the selection contains popular
works together with many of their editions -- exactly the hard cases for
edition-level matching. This choice has a cost. The overlap figures in Table 3 are higher than they would be for a random Amazon sample. 
They describe this subset and should not be generalized to all of Amazon. 
Our aim is a benchmark with enough shared editions, not a representative sample.

\begin{table}[htbp]
\centering
\begin{threeparttable}
\caption{Overview of the datasets.}
\label{tab:overview}
\footnotesize
\newcolumntype{L}[1]{>{\RaggedRight}p{#1}}
\newcolumntype{R}[1]{>{\raggedleft}p{#1}}
\newcolumntype{Y}{>{\RaggedRight\hbadness=10000}X}
\begin{tabularx}{\textwidth}{|L{1.7cm}|L{1.6cm}|l|R{1.0cm}|R{0.8cm}|Y|}
\hline
\textbf{Dataset} & \textbf{Source} & \textbf{Format} & \textbf{\# Ent.} & \textbf{\# Attr.} & \textbf{List of attributes} \\
\hline
Goodreads Best Books Ever & Kaggle / Zenodo\tnote{1} & CSV & 52,478 & 25 & bookId, title, series (MV), author, rating, description, language, isbn, genres, characters (MV), bookFormat, edition (MV), pages, publisher, publishDate, firstPublishDate (MV), awards (MV), numRatings, ratingsByStars, likedPercent, setting (MV), coverImg, bbeScore, bbeVotes, price \\
\hline
DBpedia & SPARQL endpoint\tnote{2} & XML & 59,270 & 9\tnote{4} & title, authors, publisher, release\_date, year, pages (MV), language, genres (MV), isbns (MV) \\
\hline
Amazon Reviews 2023 & Hugging Face\tnote{3} & CSV & 40,000 & 14\tnote{4} & title, authors, format, publisher, edition (MV), publication\_date, language, pages, genres, average\_rating, ratings\_count, price, isbn13, isbn10 \\
\hline
\end{tabularx}
\begin{tablenotes}\footnotesize
  \item[1] \url{https://www.kaggle.com/datasets/arnabchaki/goodreads-best-books-ever}, original: \url{https://doi.org/10.5281/zenodo.4265096}
  \item[2] \url{https://dbpedia.org/sparql}, script \texttt{dbpedia\_books.py}
  \item[3] \url{https://huggingface.co/datasets/McAuley-Lab/Amazon-Reviews-2023}, script \texttt{amazon\_subset.py}
  \item[4] Without the identifiers \texttt{uri} (DBpedia) and \texttt{amazon\_id}, \texttt{asin} (Amazon).
\end{tablenotes}
\end{threeparttable}
\end{table}

Attributes with more than 30\% missing values are marked (MV). In Goodreads
these are \texttt{series} (55.3\%), \texttt{characters} (73.8\%),
\texttt{edition} (90.6\%), \texttt{firstPublishDate} (40.6\%),
\texttt{awards} (79.8\%) and \texttt{setting} (77.9\%); in DBpedia
\texttt{pages} (38.4\%), \texttt{genres} (36.4\%) and \texttt{isbns} (39.0\%);
in Amazon only \texttt{edition} (40.1\%).

\section{Integrated Schema}

Table~\ref{tab:schema} maps the source attributes to the integrated schema.
Compared with the first proposal, \texttt{format} and \texttt{edition} were
added because they distinguish editions of the same work, and the publication
date is kept per edition instead of being reduced to a year.

\begin{table}[htbp]
\centering
\begin{threeparttable}
\caption{Attribute intersection with the integrated schema (one record per edition).}
\label{tab:schema}
\footnotesize\setlength{\tabcolsep}{4pt}
\begin{tabularx}{\textwidth}{|l|l|>{\RaggedRight\arraybackslash}X|}
\hline
\textbf{Attribute name} & \textbf{Type} & \textbf{Datasets in which the attribute is found} \\
\hline
title             & string           & GR (title), DB (title), AM (title) \\
authors           & list$<$string$>$ & GR (author), DB (authors), AM (authors) \\
genres            & list$<$string$>$ & GR (genres), DB (genres), AM (genres) \\
publication\_date & date             & GR (publishDate), DB (release\_date), AM (publication\_date) \\
publisher         & string           & GR (publisher), DB (publisher), AM (publisher) \\
page\_count       & integer          & GR (pages), DB (pages), AM (pages) \\
language          & string           & GR (language), DB (language), AM (language) \\
format            & string           & GR (bookFormat), AM (format) \\
edition           & string           & GR (edition), AM (edition) \\
rating            & float            & GR (rating), AM (average\_rating) \\
ratings\_count    & integer          & GR (numRatings), AM (ratings\_count) \\
price             & float            & GR (price), AM (price) \\
series            & string           & GR (series) \\
awards            & list$<$string$>$ & GR (awards) \\
\hline
\multicolumn{3}{|l|}{\emph{Identifier for overlap estimation and the gold standard only (not used for matching):}} \\
isbn              & string           & GR (isbn), DB (isbns), AM (isbn13, isbn10) \\
\hline
\end{tabularx}
\begin{tablenotes}[para,flushleft]
  \footnotesize GR = Goodreads, DB = DBpedia, AM = Amazon; source column names in brackets.
\end{tablenotes}
\end{threeparttable}
\end{table}

Twelve attributes occur in at least two datasets and seven in all three
(\texttt{title}, \texttt{authors}, \texttt{genres}, \texttt{publication\_date},
\texttt{publisher}, \texttt{page\_count}, \texttt{language}), which allows
fusion by voting. Representations differ: Amazon titles carry series or edition
details in brackets (``Where the Crawdads Sing (Movie Tie-in)''), Goodreads
stores contributors and roles in one string (``J.K. Rowling, Mary GrandPr\'e
(Illustrator)''), DBpedia dates range from years to full dates, and the
sources use different genre vocabularies.

\section{Overlap Estimation}

Two records describe the same edition if they share an ISBN (ISBN-10 and
ISBN-13 normalised to ISBN-13). The ISBN is used only for this estimate and as
the gold standard; the matchers in Phase~II will not use it. Since records
without an ISBN cannot be counted, the numbers in Table~\ref{tab:overlap} are
lower bounds.

\begin{table}[htbp]
\centering
\caption{Edition overlap (shared ISBN-13) and comparison with a random Amazon sample.}
\label{tab:overlap}
\footnotesize
\begin{tabular}{lr}
\toprule
\textbf{Measure} & \textbf{Editions} \\
\midrule
Amazon $\cap$ Goodreads & 2,169 \\
Amazon $\cap$ DBpedia & 1,725 \\
Goodreads $\cap$ DBpedia & 3,271 \\
In all three datasets & 539 \\
\textbf{In at least two datasets} & \textbf{6,087} \\
\midrule
Random Amazon sample (40,000, seed 42) $\cap$ Goodreads & 331 \\
\bottomrule
\end{tabular}
\end{table}

With 6,087 editions in at least two datasets, the requirement of 1,000 is met
even under this strict criterion. The total number of distinct editions is
well above 2,500, since the three sources contain 151,748 records with at most
6,087 shared editions.

To show the difficulty of edition-level matching, we also compared works
(normalised title and first-author surname). Goodreads and DBpedia share
10,619 works; for 4,759 of the 7,561 shared works with ISBNs on both sides
(63\%), the two sources list \emph{different} editions. Amazon and Goodreads
share 3,146 works (22\% in different editions), and within the Amazon subset
6,588 works appear in more than one edition (33,540 records). A matcher based
on title and author alone would therefore produce many false matches; date,
publisher, format and page count are needed to tell editions apart.
In Phase II, the matchers will therefore compare title, authors, publisher, publication date, format and page count.
Title and authors first generate candidate pairs. The other attributes then decide whether two candidates are the same edition.
We evaluate against the ISBN gold standard.


\section{Phase 1: Schema Mapping and Data Translation}

\subsection{Target Schema}
The integrated schema is defined as an XML Schema (\texttt{target.xsd}) and
serves as the common target for all three sources. Its root element
\texttt{books} contains a sequence of \texttt{book} elements, each describing
exactly one book edition, in line with our entity definition. Every book
carries a required \texttt{id} attribute that is prefixed with the name of its
source (e.g.\ \texttt{goodreads\_id\_1}). These IDs keep each record traceable
to its origin and serve as keys in the correspondences produced during
identity resolution.

The schema contains 14 attributes. Only \texttt{title} is mandatory; all other
elements are optional (\texttt{minOccurs="0"}). We chose this deliberately: no
source covers all attributes, and a stricter schema would either reject valid
records or force us to insert artificial default values that would later
distort data fusion. Missing values are therefore simply omitted.

We assigned explicit data types so that later steps can compare values without
additional parsing: \texttt{page\_count} and \texttt{ratings\_count} are
\texttt{xs:integer}, \texttt{rating} and \texttt{price} are
\texttt{xs:decimal}, and \texttt{publication\_date} is \texttt{xs:date}.
Multi-valued attributes are modelled as nested lists instead of delimited
strings: \texttt{authors}, \texttt{genres} and \texttt{awards} contain one
child element per value. This enables list-based comparators during identity
resolution and union-based fusion of genres. Compared with our proposal,
\texttt{format} and \texttt{edition} were added because they help to
distinguish editions of the same work, and the full publication date is kept
instead of only the year.

The \texttt{series} attribute remains in the target schema but is filled
only from Goodreads, the only source with a dedicated field (23,469
records, 44.7\%). In Amazon, 38.4\% of the titles contain brackets, but only
a small share (about 3\%) clearly looks like a numbered series; most
brackets describe the edition instead, such as ``Spanish Edition'' or
``Random House Large Print''. In DBpedia, 21.3\% of the titles contain
brackets, almost exclusively Wikipedia disambiguations such as ``(novel)''
(6,446) or ``(book)'' (1,857). Since the series cannot be extracted
reliably for Amazon and DBpedia, we do not use it for identity resolution
or data fusion.

\subsection{Mappings}
We implemented one Altova MapForce mapping per source. Goodreads and Amazon
are read as CSV files, DBpedia as the XML file produced by our SPARQL download
script. Each mapping writes an XML file that is valid against
\texttt{target.xsd}. Goodreads provides the richest attribute set, including
series, awards, rating, ratings count and price, but stores several values in
single strings. Amazon provides format and edition as separate fields and
lists authors and genres as strings that need to be split. DBpedia
contributes title, authors, publisher, release date, pages, language and
genres. Each record receives a source-prefixed ID generated with
\texttt{auto-number} and \texttt{concat}. Because these IDs depend on the
order of the output, the XML files must not be regenerated without re-running
identity resolution; otherwise the stored correspondences point to the wrong
records.

\subsection{User-Defined Functions}
Many cleaning steps are needed in more than one mapping. Instead of rebuilding
the same chains of built-in functions, we implemented reusable user-defined
functions (UDFs), which keep the mappings readable and ensure that the same
rule is applied everywhere (Table~\ref{tab:udf}).

\begin{table}[htbp]
\centering
\caption{User-defined functions in the MapForce mappings.}
\label{tab:udf}
\footnotesize
\begin{tabularx}{\textwidth}{@{}l l >{\RaggedRight\arraybackslash}X@{}}
\toprule
\textbf{Function} & \textbf{Used for} & \textbf{Purpose} \\
\midrule
\texttt{Clean\_Authors} & GR author & splits the author string into individual author elements \\
\texttt{Clean\_List} & GR genres, awards; AM authors & splits multi-valued strings and drops empty entries \\
\texttt{Parse\_Date} & GR publishDate & detects four date formats and converts them to \texttt{xs:date} \\
\texttt{Page\_Count} & pages & extracts the number from values such as ``1 page'' \\
\texttt{Safe\_Number} & GR price & converts only valid numbers; other values stay empty \\
\bottomrule
\end{tabularx}
\end{table}

\subsection{Date Formats}
The sources differ strongly in date representation and precision
(Table~\ref{tab:dates}).

\begin{table}[htbp]
\centering
\caption{Precision of the publication date per source.}
\label{tab:dates}
\footnotesize
\begin{tabular}{lrrr}
\toprule
 & \textbf{Goodreads} & \textbf{DBpedia} & \textbf{Amazon} \\
\midrule
Full date               & 84.1\% & 22.2\% & 95.1\% \\
Year only               &  8.3\% & 58.0\% &  0.0\% \\
Other (e.g.\ month--year) & 5.9\% & 10.5\% & 0.0\% \\
Empty                   &  1.7\% &  9.2\% &  4.9\% \\
\bottomrule
\end{tabular}
\end{table}

Goodreads uses four formats: US dates with a two-digit year (04/27/10), with a
four-digit year (04/27/2010), written dates with an ordinal suffix (April 27th
2010) and bare years (2003). Amazon uses written dates such as January 1,
2019. DBpedia combines several infobox properties, so 58\% of its values
contain only a year, followed by ISO dates (1982-04-15) and month--year values
(July 1941).

\texttt{Parse\_Date} first identifies the Goodreads format with regular
expressions (\texttt{match-pattern}) and then applies a dedicated conversion
in a cascade of \texttt{if-else} components. Two-digit years are completed
with the century (above 26 $\rightarrow$ 19xx, otherwise 20xx), four-digit US
dates are parsed with the pattern \texttt{[M01]/[D01]/[Y0001]}, ordinal
suffixes are removed with \texttt{substring-before}, \texttt{substring-after}
and \texttt{string-length} before parsing with \texttt{[MNn] [D1] [Y0001]},
and bare years are completed to 1~January. Values matching none of the
patterns are left empty, so no invalid dates are produced. For Amazon, the
date is matched with a regular expression and converted with
\texttt{parse-date}. % TODO: DBpedia

Completing year-only values to 1~January makes these dates look more precise
than they are. In Amazon, 17.6\% of the dates fall on 1~January, which
suggests that the source itself fills unknown days with this default.

\subsection{Numeric Values and Further Problems}
A single invalid value in an \texttt{xs:integer} or \texttt{xs:decimal} field
makes MapForce abort the transformation. We therefore read numeric fields as
strings and convert them with UDFs. \texttt{Safe\_Number} converts the
Goodreads price only if it is a valid number; malformed values such as
1.189.88 (12 records) and 14,365 empty values are omitted, leaving 38,101
valid prices. \texttt{Page\_Count} extracts the number from values such as
``1 page'' (23 records), so that 50,131 Goodreads records carry a page count.
Splitting Amazon genres produced empty tokens, which we removed with a
filter. Finally, some Goodreads award strings contain escaped HTML entities
(\texttt{\&amp;}) that already exist in the source CSV; since they do not
affect matching, we left them unchanged.

\end{document}
